\scanpage{214}
\mspage{163}

soit un isom. local, ou ce qui revient au $=$, que \deleted{\unsure{ce soit un}} $X'$ soit \uline{\unsure{fibrée|fibré}} sur $X$ à fibres discrètes :

\[
\varphi_X\colon \unsure{X|{}'X}^\wedge
\overset{\approx}{\mathrel{\unsure{-|\longrightarrow}}}
\underbrace{\underline{\unsure{\mathrm{Faib}|\mathrm{Fib}}\mathrm{.disc.}}\,X}_{\substack{\text{catégories fibrées}\\\text{à fibres discrètes}}} ,
\]

Notons que toute flèche $X'\unsure{-|\longrightarrow}X''$ dans $\underline{\mathrm{Fib}}\,\underline{\mathrm{disc}}\,X$ \deleted{\unsure{qui}} est un isom. local (tautologique sur

\diagram{\begin{tikzcd}[ampersand replacement=\&,row sep=small,column sep=small]
X'_{/x'}\arrow[rr,dash]\arrow[dr,dash,"\mathrm{iso}"'] \& \& X''_{/x''}\arrow[dl,dash,"\mathrm{iso}"]\\
\& X_{/x} \&
\end{tikzcd}}

), ce qui montre que si on a un diagramme

\diagram{\begin{tikzcd}[ampersand replacement=\&,row sep=small,column sep=small]
X_1\arrow[dr]\& \& X_2\arrow[dl]\\
\& X_0\&
\end{tikzcd}}

dans $\underline{\mathrm{Fib}}.\underline{\mathrm{disc}}.\,X$, son produit fibré $X_3=\unsure{\alpha_1|X_1}\times_{X_0}X_2$ dans Cat est tel que $X_3\unsure{-|\longrightarrow}X_2$, est un iso. local, donc aussi $X_3\longrightarrow X$. C'est-à-dire $\underline{\mathrm{Fib.disc}}\,X$ est stable par produits fibrés. Ceci implique que les \deleted{foncteur $\varphi_X$ transforme produits fib} \inserted{produits fibrés existent dans} $\underline{\mathrm{Fib\,disc}}\,X$, et que ce sont les produits fibrés dans Cat, ou dans $\mathrm{Cat}_{/\unsure{x|X}}$.
% archive sheet number: 214
